
\partintro{这部分是力学。}
\partimage{\includegraphics[height=\paperheight]{starry.jpg}}
\part[直线运动]{运动学I：直线运动}


\chaptersaying{不了解运动，就不了解自然。——亚里士多德}

%\chapterimage{\includegraphics[width=\paperwidth]{Newton.jpg}}
\chapter[运动的描述]{运动的描述}

\addmarginparintoc{\cutimage{1.5\linewidth}{-3.0mm}{0}{iris.jpg}}



	\begin{Block}%[测试文字]
		在我们周围，到处可以看到物体在运动：汽车在公路上行驶，巨轮在海上航行，飞机在天空中飞行，树叶在摇动，鸟儿在飞翔……就连我们脚下的地球，也在不停地自转、公转。
		
		物体的空间位置随时间的变化，是自然界中最简单、最基本的运动形态，叫作机械运动（mechanical motion）。在物理学中，研究物体做机械运动规律的分支叫作力学（mechanics）。人们在力学的研究中，不仅认识了物体做机械运动的规律，而且还创立了科学研究的基本方法。
		
		这一章，我们学习机械运动的描述。
	\end{Block}
	
	
%	\clearpage
	\begin{Proof}
		\zhlipsum
	\end{Proof}
	

\section{质点与参考系\quad 时间与位移}

	\begin{Point*}
		生活中随处可见运动的物体，例如玩耍的孩童、行驶的汽车、翱翔的雄鹰……对于这些运动的物体，我们如何准确地描述它们的运动呢？
		
		要讨论物体位置随时间的变化，就要涉及位置、时间等概念。如果要准确地描述一辆行驶在北京长安街上的汽车所处的位置，你认为应该采用什么方法？你对时间是怎样认识的？
		\tcbsubtitle{在本节，你将学到}
		\begin{UnnumberedItem}[2][itemindent=2.0em]
			\item 质点以及理想化模型；
			\item 参考系的概念和分类；
			\item 时刻与时间间隔的区别；
			\item 位移的概念；
			\item 矢量和标量。
		\end{UnnumberedItem}
	\end{Point*}


\begin{paracol}{2}
	\NoteArea[0.6cm]

	\subsection{物体和质点}
	
	
	\subsubsection{质点}
	任何物体都有一定的大小和形状，物体各部分的运动情况一般说来并不一样。不过，我们有时关注物体各部分的运动，有时关注物体整体的运动。所以，如何描述物体的运动要看我们关注的问题是什么。
	
	有些物体上各点的运动差异虽然不小，但如果我们{\ep 不研究各点的运动差异，而只关注物体整体的运动时}，就可以把物体看作一个点。
	
	还有一种情况，虽然不能忽略物体的大小和形状，但{\ep 物体上各点的运动情况完全相同，仅从描述运动的角度看，物体上任意一点的运动完全能反映整个物体的运动}，此时，整个物体的运动也可以简化为一个点的运动。
	
	在上述情况下，可以{\ep 忽略物体的大小和形状}，把它简化为一个{\ep 具有质量的点}，这样的点叫作{\ep 质点}（mass point）。


\end{paracol}

\begin{Check}
	\begin{NumberedItem}{1}
		\item 能将物体看作质点的条件：\underline{\hspace*{7.0cm}}。
		
		\item 将下列问题按照“是否能将研究对象看作质点”进行分类：
		
		\hspace*{1.5em}
		\begin{minipage}[t]{\linewidth-1.5em}
			\cirone~地球自转；
			\cirtwo~月球公转；
			\cirthree~测量高铁从成都行驶到沈阳的时间；
			\cirfour~测量一辆列车从开始进入隧道到完全离开隧道的时间；
			\cirfive~研究舞蹈演员的动作；
			\cirsix~研究橡胶的拉伸变形；
			\cirseven~研究走钢丝表演者能否稳定地站在钢丝上；
			
			可以将研究对象看作质点：\underline{\hspace*{7.0cm}}；
			
			不能将研究对象看作质点：\underline{\hspace*{7.0cm}}。
		\end{minipage}
		
	
	\end{NumberedItem}
\end{Check}



\begin{paracol}{2}

在物理学中，突出问题的主要因素，忽略次要因素，建立理想化的物理模型，并将其作为研究对象，是经常采用的一种科学研究方法。{\ep 质点}这一{\ep 理想化模型}就是这种方法的具体应用。

\subsection{参考系和坐标系}

\subsubsection{参考系}

自然界的一切物体都处于永恒的运动中，绝对静止的物体是不存在的。就此意义而言，我们说运动是绝对的。但是，描述某个物体的位置随时间的变化，却又总是相对于其他物体而言的。这便是运动的相对性。

要描述一个物体的运动，首先要选定某个其他物体作为参考，观察物体的位置相对于这个“其他物体”是否随时间变化，以及怎样变化。{\ep 这种用来作为参考的物体叫作参考系}（reference frame）。

{\ep 通常情况下，在讨论地面上物体的运动时，都以地面为参考系。}


\subsubsection{坐标系}

为了定量描述物体相对于参考系的运动，需要建立坐标系。坐标系包含数字和轴，建立在参考系上，相对于参考系静止。


\end{paracol}


\begin{Check*}
	一个小球从距地面5\,m高处落下，被地面弹回，在距地面1\,m高处被接住，坐标原点定在释放点正下方2\,m处，竖直向下为坐标轴正方向建立坐标系如图所示。则小球的释放点、落地点、接住点的位置坐标分别是\,\bkt
	

	\choice{$-$2\,m、3\,m、2\,m}{$-$2\,m、2\,m、1\,m}{$-$2\,m、$-$2\,m、$-$1\,m}{4\,m、0、1\,m}
\end{Check*}


\begin{paracol}{2}
	
	
	\subsection{时刻与时间间隔}
	
	\subsubsection{时刻与时间}
	
	
	
	
	
	\subsection{位置与位移}
	
	\subsubsection{路程与位移}
	
	\subsubsection{标量和矢量}
	
	\subsubsection{直线运动的位移}
	
	\subsubsection{位移-时间图像}
	
	\subsection{}
	
	
\end{paracol}



\Thinking[Read and Think]

\Exercise

\NewDocumentCommand{\mytblrs}{m m +m}{%
\noindent
\begin{tblr}{
width=\linewidth,
rows={#1},
colspec={#2},
row{1}={bg=ColorB!7.5,font=\bfseries},
hlines={gray},
hline{1,Z}={2pt,ColorA},
hline{2}={1pt,ColorA},
vline{2-Y}={gray}
}
#3
\end{tblr}}


\mytblrs
{rowsep=1ex}    % 行高
{X[1,c,m]|X[6,c,m]|X[2,c,m]}   % 列宽
{名称 & 定义 & 表示方法 \\
集合 & 由若干确定的、互不相同的对象所组成的整体 & $A,B,C,\cdots$ \\
元素 & 组成集合的每个对象 & $a,b,c\cdots$ \\}


\chapter{短标题}

{\FontSizeSet{20pt}测试文字测试文字测试文字测试文字测试文字测试文字测试文字测试文字测试文字测试文字测试文字测试文字测试文字测试文字}
\begin{paracol}{2}
%	添加脚注\footnote{脚注}
	\marginnote{测试文字测试文字测试文字测试文字测试文字测试文字测试文字测试文字测试文字测试文字测试文字测试文字测试文字测试文字}
	
%	\makeatletter
%	\show\output
%	\show\@outputpage
%	\makeatother
	
	\zhlipsum[1]\footnote{测试脚注位置}
	
	\zhlipsum[2]
	
	\begin{Warning}
		\lipsum
		\tcbsubtitle{测试文字}
		\tcblower
		\tcbsubtitle{测试文字}
		cswz
	\end{Warning}
	
	\zhlipsum[3]
\end{paracol}


\begin{Variant}
这是一道例题这是一道例题这是一道例题这是一道例题这是一道例题这是一道例题这是一道例题这是一道例题这是一道例题这是一道例题这是一道例题这是一道例题这是一道例题这是一道例题这是一道例题这是一道例题这是一道例题这是一道例题这是一道例题这是一道例题这是一道例题这是一道例题这是一道例题这是一道例题这是一道例题这是一道例题这是一道例题这是一道例题这是一道例题这是一道例题这是一道例题这是一道例题这是一道例题这是一道例题这是一道例题这是一道例题这是一道例题这是一道例题这是一道例题这是一道例题这是一道例题这是一道例题这是一道例题这是一道例题

\begin{NumberedItem}{1}
	\item 能将物体看作质点的条件：\underline{\hspace*{7.0cm}}。
	
	\item 将下列问题按照“是否能将研究对象看作质点”进行分类：
		\begin{NumberedItem}{2}
			\item 能将物体看作质点的条件：\underline{\hspace*{7.0cm}}。
			
			\item 将下列问题按照“是否能将研究对象看作质点”进行分类：
		\end{NumberedItem}
\end{NumberedItem}

\end{Variant}


\begin{NumberedItem}{1}
	\item 能将物体看作质点的条件：\underline{\hspace*{7.0cm}}。
	
	\item 将下列问题按照“是否能将研究对象看作质点”进行分类：
\end{NumberedItem}



\Formula{\int_\Omega\left(x^2+x\right){\rm d}x}



\begin{enumerate}
	\item 能将物体看作质点的条件：\underline{\hspace*{7.0cm}}。
	
	\item 将下列问题按照“是否能将研究对象看作质点”进行分类：
		\begin{enumerate}
			\item 能将物体看作质点的条件：\underline{\hspace*{7.0cm}}。
			
			\item 将下列问题按照“是否能将研究对象看作质点”进行分类：
		\end{enumerate}
		
\end{enumerate}



\clearpage

\begin{Example}
{\bfseries\color{ColorA}(2019年4月浙江物理选考)}~如图所示，一根粗糙的水平横杆上套有$ A $、$ B $两个轻环，系在两环上的登场细绳拴住的书本处于静止状态，现将两环距离变小后书本仍处于静止状态，则~{\bfseries\color{ColorA}B}

\bigskip
\noindent%
\begin{minipage}{0.7\linewidth}
\choice{杆对$ A $环的支持力变大}{$ B $环对杆的摩擦力变小}{杆对$ A $环的力不变}{与$ B $环相连的细绳对书本的拉力变大与$ B $环相连的细绳对书本的拉力变大}
\end{minipage}
\hspace{5.0mm}
\begin{minipage}{0.25\linewidth}
\includegraphics[width=\linewidth]{430.pdf}
\end{minipage}

\end{Example}


\begin{Answer}
	测试文字
\end{Answer}







%\chapterimage{\includegraphics[height=\paperheight]{poppies.jpg}}
\chapter{附录1}
%\AddGutter



\begin{multicols}{2}
%\raggedcolumns%
\zhlipsum\zhlipsum
\end{multicols}

\clearpage

\begin{multicols}{2}
\zhlipsum
\end{multicols}




\section{附录测试}

\begin{Appendix}

%\chapterimage{\includegraphics[height=\paperheight]{iris1.jpg}}
\chapter{附录2}


\begin{paracol}{2}
\lipsum
%\switchcolumn
%\sloppy
%\lipsum[1-2]
\end{paracol}


\section{附录测试}







	

\section{附录测试}



\begin{Appendix*}
\section{附录测试}

\makeatletter
\@SubsectionStyle{测试文字测试文字测试文字测试文字测试文字测试文字测试文字}{}{1}
\makeatother




\subsection{测试文字节层次附录附录}

\subsubsection{看目录}

\subsubsection{看目录}

\subsubsection*{看目录}


\section{测试}

\end{Appendix*}

\section{附录测试}

	
\section{附录测试}
	
\end{Appendix}


\begin{Project}
	\section{附录测试}
	
	\subsubsection{看目录看目录}
	
	\subsubsection{看目录}
	
	\subsubsection*{看目录}
	
	
	\section{测试}
	
\end{Project}


\section{测量物体，包含长度测量}

\subsubsection{看目录}

\subsubsection{看目录}

\subsubsection*{看目录}



\begin{Quiz}
	\section{附录测试}
	
	\subsection{测试文字}
	
	\subsubsection{看目录}
	
	\subsubsection{看目录}
	
	\subsubsection*{看目录}
	
	\subsection{测试文字}
	
	\subsection{测试文字}
	
	
	\section{测试}
	

	
\end{Quiz}

\partimage{\includegraphics[height=\paperheight]{starry.jpg}}
\part{丝柏树}



\chapterimage{\includegraphics[height=\paperheight]{cypresses.jpg}}
\chapter{彩色计时器彩色计时器彩色计时器彩色计时器彩色计时}

文森特·威廉·梵高（Vincent Willem van Gogh，1853年3月30日—1890年7月29日），荷兰后印象派画家。代表作有《星月夜》、自画像系列、向日葵系列等。

梵高出生于荷兰乡村津德尔特的一个新教牧师家庭，早年的他做过职员和商行经纪人，还当过矿区的传教士，最后他投身于绘画。他早期画风写实，受到荷兰传统绘画及法国写实主义画派的影响。1886年，他来到巴黎，结识印象派和新印象派画家，并接触到日本浮世绘的作品，视野的扩展使其画风巨变。1888年，来到法国南部小镇阿尔，创作《阿尔的吊桥》；同年与画家保罗·高更交往，但由于二人性格的冲突和观念的分歧，合作很快便告失败。此后，梵高的疯病（有人记载是“癫痫病”）时常发作，但神志清醒时他仍然坚持作画。1889年创作《星月夜》。1890年7月，梵高在精神错乱中开枪自杀（一说，两个年轻人不小心走火开枪击中），年仅37岁。

梵高是表现主义的先驱，并深深影响了二十世纪艺术，尤其是野兽派与德国表现主义。\footnote{本章封面画为荷兰画家文森特·梵高的《柏树》。}





\section{测量物体，包含长度测量}

\footnote{测试文字}


\noindent{\Huge\score{4.4}}


\noindent\score{4.4}测试文字测试文字测试文字测试文字测试文字测试文字测试文字测试文字测试文字测试文字测试文字测试文字测试文字测试文字测试文字测试文字测试文字测试文字测试文字测试文字测试文字测试文字测试文字测试文字测试文字测试文字测试文字测试文字测试文字测试文字测试文字测试文字测试文字测试文字测试文字测试文字测试文字测试文字测试文字测试文字测试文字测试文字测试文字测试文字测试文字测试文字测试文字



\addmarginparintoc{\cutimage{1.5\linewidth}{0}{0}{olive.jpg}}[1.0cm]



\addmarginparintoc{\cutimage{1.5\linewidth}{0}{0}{harvest.jpg}}[1.0cm]


\addmarginparintoc{\cutimage{\linewidth}{0}{-5.0mm}{cypresses.jpg}}[1.0cm]

%\protect
\begin{Project}

\section{测量物体，包含长度测量}

\begin{Point*}
	学习完这一单元，你应当能够: 识别国际单位制(SI)中的基本量; 正确称呼国际单位制中最常用的词头; 利用链环变换法改变单位(这里是对于长度、面积和体积); 说明米是如何用真空中光速来定义的。
	\tcblower
	\tcbsubtitle{本节重点\footnote{脚注}}
	物理学是以物理量的测量为基础的。一些物理量被选作基本物理量; 每一种基本物理量都已经通过一种标准来定义并给定测量的单位。另一些物理量是用这些基本物理量和它们的标准及单位来定义。

\end{Point*}

测试文字测试文字测试文字测试文字测试文字测试文字测试文字测试文字测试文字测试文字测试文字测试文字测试文字测试文字测试文字测试文字测试文字测试文字测试文字测试文字测试文字测试文字测试文字测试文字测试文字测试文字




\begin{paracol}{2}
	\NoteArea[2.0em]
	
	
%	\tcbox[boxsep=0mm,fontupper=\BoxTitleFont,colupper=ColorA,colframe=ColorB!7.5,width=\linewidth]{\faCaretLeft~\notetitle~\faCaretRight}
	
	\subsection{什么是物理学？}
	
	你这样也太没有 \leaders \hrule\hskip 2cm plus 0.5cm minus 0.5cm 了。
	
	我一定会让你 \leaders \hrule height
	-2.6pt depth 3pt \hskip 2cm 满意的。
	
	
	科学和工程学建立在测量和比较的基础上。因此，我们就需要关于如何测量和比较物体的规则，并且我们还需要通过实验来建立这种测量和比较的单位。物理学(以及工程学)的一个目的就是设计和实施这些实验\DeepThink*{测试文字测试文字测试文字测试文字测试文字测试文字}。
	
%	\marginnote{测试文字测试文字测试文字测试文字测试文字}
	
	例如，物理学家们尽力开发极其准确的时钟从而使任何时间或时间间隔都可以精确地确定和比较。你们可能要问，这样的准确度是否真正需要，或者是否值得花这么大的精力。这里给出一个是否值得的例子：如果没有极其准确的时钟的话，那么现在全世界范围的航行必不可少的全球定位系统(GPS)就会变得毫无意义\DeepThink*{测试文字测试文字测试文字测试文字测试文字测试文字}。
	
%	\NoteArea
	\subsection{测量物体\protect\footnote{看看我}}
%	\NoteArea
	
	\subsubsection{测试文字}
	
	我们通过学习怎样测量物理学中的量来学习物理学。这些物理量包括长度、时间、质量、温度、压强和电流。
	
	我们通过和一个标准相比较，用它们各自的单位来量度各个
	
	\begin{Lemma}
		测试文字
	\end{Lemma}
	
	
	
\end{paracol}

\subsection{力的三要素}


\subsubsection*{测试文字}

力是有大小的．我们在初中学过，力的大小可以用测力计来测量．在国际单位制中力的单位是\textbf{牛顿}，简称牛，国际符号是N，日常生活和生产中常用的力的单位是千克力，牛顿和千克力的关系是：1千克力$=9.8$牛．

力不但有大小，而且有方向．物体受到的重力是竖直向下的，物体在液体中受到的浮力是竖直向上的，力的方向不同，它的作用效果也不同．用力拉弹簧，弹簧就伸长；用反方向的力压弹簧，弹簧就缩短．作用在运动物体上的力，如果方向与运动方向相同，将加快物体的运动；如果方向与运动方向相反，将阻碍物体的运动．可见，要把一个力完全表达出来，除了说明力的大小外，还要指明力的方向．


为了直观地说明力的作用，常常用一根带箭头的线段来表示力．线段是按一定比
例（标度）画出的，它的长短表示力的大小，它的指向表示力的方向，箭头或箭尾表示力的作用点，箭头所沿的直线叫做力的作用线．这种表示力的方法，叫做\textbf{力的图示}．图1.1中表示的是作用在小车上的$100$N的力．

\subsection{力的分类\protect\footnote{测试文字}}

我们从初中开始学习物理以来，见过的力的名称已经相当多了．各种力可以用不同的方法来分类．一种是根据力的性质来分类的，如重力、弹力、摩擦力、分子力、电磁力等；另一种是根据力的效果来分类的，如拉力、压力、支持力、动力、阻力等等．拉力、压力、支持力实际上都是弹力，只是效果不同．不论是什么性质的力，只要效果是加快物体的运动，就可以叫它为动力；效果是阻碍物体的运动，就可以叫它为阻力．今后我们还会遇到根据效果来命名的力的名称．

从力的性质来看，力学中经常遇到的有重力、弹力、摩擦力．下面几节就分别介绍这三种力．


\begin{Example}
	\lipsum[1][1-7]
	\begin{flalign}
		\Psi=\int_\Omega\left[f\left(\phi\right)+\kappa\left|\nabla\phi\right|^2\right]{\rm d}V
	\end{flalign}
\end{Example}

\begin{Answer*}
	测试文字测试文字\footnote{看看我}
\end{Answer*}


\subsection{测试文字}
\subsection*{测试文字}
\zhlipsum[1]


\makeatletter



\makeatother

\begin{multicols}{2}
\Basic

\Complex
\end{multicols}

\zhlipsum[1]

















\clearpage

\Improve
\begin{QuestionItem}[2]{1}
	\item 关于行星运动的规律，下列说法符合史实的是
	\choice{开普勒在牛顿定律的基础上，导出了行星运动的规律}{开普勒在天文观测数据的基础上，总结出了行星运动的规律}{开普勒总结出了行星运动的规律，找出了行星按照这些规律运动的原因}{开普勒总结出了行星运动的规律，发现了万有引力定律}
	\item 火星和木星沿各自的椭圆轨道绕太 阳运行，根据开普勒行星运动定律可知
	\choice{太阳位于木星运行轨道的中心}{火星和木星绕太阳运行速度的大小始终相等}{火星与木星公转周期之比的平方等于它们轨道半长轴之比的立方}{相同时间内，火星与太阳连线扫过的面积等于木星与太阳连线扫过的面积}
	\item 为了探测引力波，“天琴计划” 预计发射地球卫星P，其轨道半径约为地球半径的16倍；另一地球卫星Q的轨道半径约为地球半径的4倍。P与Q的周期之比约为
	\choice{2∶1}{4∶1}{8∶1}{16∶1}
	\item 测试文字测试文字测试文字测试文字测试文字测试文字测试文字测试文字测试文字测试文字测试文字测试文字测试文字测试文字测试文字测试文字测试文字测试文字测试文字测试文字测试文字测试文字测试文字测试文字测试文字测试文字测试文字测试文字测试文字测试文字测试文字测试文字测试文字测试文字测试文字测试文字测试文字测试文字测试文字测试文字测试文字测试文字测试文字测试文字测试文字测试文字测试文字测试文字
	\item 测试文字测试文字测试文字测试文字测试文字测试文字测试文字测试文字测试文字测试文字测试文字测试文字测试文字测试文字测试文字测试文字测试文字测试文字测试文字测试文字测试文字测试文字测试文字测试文字测试文字测试文字测试文字测试文字测试文字测试文字测试文字测试文字测试文字测试文字测试文字测试文字测试文字测试文字测试文字测试文字测试文字测试文字测试文字测试文字测试文字测试文字测试文字测试文字
	\item 测试文字测试文字测试文字测试文字测试文字测试文字测试文字测试文字测试文字测试文字测试文字测试文字测试文字测试文字测试文字测试文字测试文字测试文字测试文字测试文字测试文字测试文字测试文字测试文字测试文字测试文字测试文字测试文字测试文字测试文字测试文字测试文字测试文字测试文字测试文字测试文字测试文字测试文字测试文字测试文字测试文字测试文字测试文字测试文字测试文字测试文字测试文字测试文字
	\item 测试文字测试文字测试文字测试文字测试文字测试文字测试文字测试文字测试文字测试文字测试文字测试文字测试文字测试文字测试文字测试文字测试文字测试文字测试文字测试文字测试文字测试文字测试文字测试文字测试文字测试文字测试文字测试文字测试文字测试文字测试文字测试文字测试文字测试文字测试文字测试文字测试文字测试文字测试文字测试文字测试文字测试文字测试文字测试文字测试文字测试文字测试文字测试文字
	\item 测试文字测试文字测试文字测试文字测试文字测试文字测试文字测试文字测试文字测试文字测试文字测试文字测试文字测试文字测试文字测试文字测试文字测试文字测试文字测试文字测试文字测试文字测试文字测试文字测试文字测试文字测试文字测试文字测试文字测试文字测试文字测试文字测试文字测试文字测试文字测试文字测试文字测试文字测试文字测试文字测试文字测试文字测试文字测试文字测试文字测试文字测试文字测试文字
	\item 测试文字测试文字测试文字测试文字测试文字测试文字测试文字测试文字测试文字测试文字测试文字测试文字测试文字测试文字测试文字测试文字测试文字测试文字测试文字测试文字测试文字测试文字测试文字测试文字测试文字测试文字测试文字测试文字测试文字测试文字测试文字测试文字测试文字测试文字测试文字测试文字测试文字测试文字测试文字测试文字测试文字测试文字测试文字测试文字测试文字测试文字测试文字测试文字
	\item 测试文字测试文字测试文字测试文字测试文字测试文字测试文字测试文字测试文字测试文字测试文字测试文字测试文字测试文字测试文字测试文字测试文字测试文字测试文字测试文字测试文字测试文字测试文字测试文字测试文字测试文字测试文字测试文字测试文字测试文字测试文字测试文字测试文字测试文字测试文字测试文字测试文字测试文字测试文字测试文字测试文字测试文字测试文字测试文字测试文字测试文字测试文字测试文字
	\item 测试文字测试文字测试文字测试文字测试文字测试文字测试文字测试文字测试文字测试文字测试文字测试文字测试文字测试文字测试文字测试文字测试文字测试文字测试文字测试文字测试文字测试文字测试文字测试文字测试文字测试文字测试文字测试文字测试文字测试文字测试文字测试文字测试文字测试文字测试文字测试文字测试文字测试文字测试文字测试文字测试文字测试文字测试文字测试文字测试文字测试文字测试文字测试文字
	\item 测试文字测试文字测试文字测试文字测试文字测试文字测试文字测试文字测试文字测试文字测试文字测试文字测试文字测试文字测试文字测试文字测试文字测试文字测试文字测试文字测试文字测试文字测试文字测试文字测试文字测试文字测试文字测试文字测试文字测试文字测试文字测试文字测试文字测试文字测试文字测试文字测试文字测试文字测试文字测试文字测试文字测试文字测试文字测试文字测试文字测试文字测试文字测试文字
	\item 测试文字测试文字测试文字测试文字测试文字测试文字测试文字测试文字测试文字测试文字测试文字测试文字测试文字测试文字测试文字测试文字测试文字测试文字测试文字测试文字测试文字测试文字测试文字测试文字测试文字测试文字测试文字测试文字测试文字测试文字测试文字测试文字测试文字测试文字测试文字测试文字测试文字测试文字测试文字测试文字测试文字测试文字测试文字测试文字测试文字测试文字测试文字测试文字
	\item 测试文字测试文字测试文字测试文字测试文字测试文字测试文字测试文字测试文字测试文字测试文字测试文字测试文字测试文字测试文字测试文字测试文字测试文字测试文字测试文字测试文字测试文字测试文字测试文字测试文字测试文字测试文字测试文字测试文字测试文字测试文字测试文字测试文字测试文字测试文字测试文字测试文字测试文字测试文字测试文字测试文字测试文字测试文字测试文字测试文字测试文字测试文字测试文字
	\item 测试文字测试文字测试文字测试文字测试文字测试文字测试文字测试文字测试文字测试文字测试文字测试文字测试文字测试文字测试文字测试文字测试文字测试文字测试文字测试文字测试文字测试文字测试文字测试文字测试文字测试文字测试文字测试文字测试文字测试文字测试文字测试文字测试文字测试文字测试文字测试文字测试文字测试文字测试文字测试文字测试文字测试文字测试文字测试文字测试文字测试文字测试文字测试文字
	\item 测试文字测试文字测试文字测试文字测试文字测试文字测试文字测试文字测试文字测试文字测试文字测试文字测试文字测试文字测试文字测试文字测试文字测试文字测试文字测试文字测试文字测试文字测试文字测试文字测试文字测试文字测试文字测试文字测试文字测试文字测试文字测试文字测试文字测试文字测试文字测试文字测试文字测试文字测试文字测试文字测试文字测试文字测试文字测试文字测试文字测试文字测试文字测试文字
\end{QuestionItem}


\Thinking
	\begin{QuestionItem}{1}
		\item 测试文字测试文字测试文字测试文字测试文字测试文字测试文字测试文字测试文字测试文字测试文字测试文字测试文字测试文字测试文字测试文字测试文字测试文字测试文字测试文字测试文字测试文字测试文字测试文字测试文字测试文字测试文字测试文字测试文字测试文字测试文字测试文字测试文字测试文字测试文字测试文字测试文字测试文字测试文字测试文字测试文字测试文字测试文字测试文字测试文字测试文字测试文字测试文字
			\begin{QuestionItem}{2}
				\item 测试文字测试文字测试文字测试文字测试文字测试文字测试文字测试文字测试文字测试文字测试文字测试文字测试文字测试文字测试文字
				\item 测试文字测试文字测试文字测试文字测试文字测试文字测试文字测试文字测试文字测试文字测试文字测试文字测试文字测试文字测试文字
			\end{QuestionItem}
		\item 测试文字测试文字测试文字测试文字测试文字测试文字测试文字测试文字测试文字测试文字测试文字测试文字测试文字测试文字测试文字测试文字测试文字测试文字测试文字测试文字测试文字测试文字测试文字测试文字测试文字测试文字测试文字测试文字测试文字测试文字测试文字测试文字测试文字测试文字测试文字测试文字测试文字测试文字测试文字测试文字测试文字测试文字测试文字测试文字测试文字测试文字测试文字测试文字
	\end{QuestionItem}
	
测试文字测试文字测试文字测试文字测试文字测试文字测试文字测试文字测试文字测试文字测试文字测试文字测试文字测试文字测试文字测试文字测试文字测试文字测试文字测试文字测试文字测试文字测试文字测试文字测试文字测试文字测试文字测试文字测试文字测试文字测试文字测试文字测试文字测试文字测试文字测试文字测试文字测试文字测试文字测试文字测试文字测试文字测试文字测试文字测试文字测试文字测试文字测试文字



\begin{Definition}[测试文字\footnote{测试脚注测试脚注测试脚注测试脚注测试脚注测试脚注测试脚注测试脚注测试脚注测试脚注测试脚注测试脚注测试脚注测试脚注测试脚注测试脚注测试脚注}]
	测试文字\footnote{脚注}
	\begin{flalign}
		\Psi=\int_\Omega\left[f\left(\phi\right)+\kappa\left|\nabla\phi\right|^2\right]{\rm d}V
	\end{flalign}
\end{Definition}

\begin{Lemma}
	测试文字
\end{Lemma}

%\flushbottom
\begin{Theorem}[测试文字测试文字测试文字测试文字测试文字测试文字测试文字测试文字测试文字测试文字测试文字测试文字测试文字]
	测试文字
\end{Theorem}

\begin{Axiom}[测试文字][测试文字]
	测试文字
\end{Axiom}


\begin{Proposition}
	测试文字
\end{Proposition}


\begin{Corollary}
	\lipsum[1]
\end{Corollary}

\begin{Lemma}[测试文字\footnote{脚注}]
	测试文字\footnote{看看我}
\end{Lemma}


\begin{Lemma*}[测试文字测试文字]
	测试文字测试文字测试文字测试文字测试文字测试文字测试文字测试文字测试文字测试文字
	
	测试文字测试文字测试文字测试文字测试文字测试文字测试文字测试文字测试文字测试文字
	\zhlipsum
	测试文字测试文字测试文字测试文字测试文字测试文字测试文字测试文字测试文字测试文字测试文字测试文字
	\tcblower
	\tcbsubtitle{测试文字}
	\zhlipsum[2]
	\tcbsubtitle{副标题}
	\zhlipsum[2]\footnote{看看我}
\end{Lemma*}



\begin{linenumbers}
\lipsum[1]
\end{linenumbers}


\begin{Comprehension}[测试文字\footnote{脚注}]
	\lipsum[1]
\end{Comprehension}



\clearpage

\lipsum[1]


\LongHeading{测试文字测试文字测试文字测试文字测试文字}


\begin{multicols}{2}
\LongHeading{测试文字\footnote{看看我}}
测试文字测试文字测试文字测试文字测试文字测试文字测试文字测试文字测试文字测试文字测试文字测试文字测试文字测试文字测试文字
\LongHeading{测试文字}
测试文字测试文字测试文字测试文字测试文字测试文字测试文字测试文字测试文字测试文字测试文字测试文字测试文字测试文字测试文字
\end{multicols}

\clearpage


\Listening

\Reading

\Writing

\Speaking

\AdjustiveHeading{测试文字\footnote{看看我}}

\lipsum[1][1-5]


\begin{UnnumberedItem}[3]

\item 测试文字测试文字
\item test
\item test
\item test
\item test
\item test
\item test
\item test
\item test
\end{UnnumberedItem}

\begin{NumberedItem}{1}
	\item 测试文字测试文字测试文字测试文字测试文字测试文字测试文字测试文字测试文字测试文字测试文字测试文字测试文字测试文字测试文字测试文字测试文字测试文字测试文字测试文字测试文字测试文字测试文字测试文字测试文字测试文字测试文字测试文字测试文字
	\item test
	\item testqjply\footnote{看看我}
\end{NumberedItem}



\begin{Proof}
	\lipsum[1]\footnote{看看我}
	\tcblower
	\lipsum[1]
	\tcbsubtitle{测试文字}
	\lipsum[2]
\end{Proof}

\begin{Block}%[测试文字]
	\lipsum[1]
	\tcblower
	\footnote{测试文字}\lipsum
\end{Block}


\begin{Check}
	\lipsum[1-2]
	\tcbline
	\lipsum[1-2]
\end{Check}

\makeatletter


\begin{Warning}
	\lipsum[2]
	\tcbsubtitle{测试文字}
	\tcblower
	\tcbsubtitle{测试文字}
	cswz
\end{Warning}
\makeatother



\begin{Comprehension}[测试文字\footnote{脚注}]
	\lipsum
		
	\lipsum
	\begin{NumberedItem}{1}
		\item 测试文字
	\end{NumberedItem}
\end{Comprehension}



\Remark{\zhlipsum[1]%
\tcblower%
\zhlipsum[2]}









%\begin{PythonBox}[神经网络]
%import pytorch as pt
%import numpy as np
%import pandas as pd
%import pytorch as pt
%import numpy as np
%import pandas as pd
%import pytorch as pt
%import numpy as np
%import pytorch as pt
%import numpy as np
%import pandas as pd
%import pytorch as pt
%import numpy as np
%import pandas as pd
%import pandas as pd
%import pytorch as pt
%import numpy as np
%import pandas as pd
%import pytorch as pt
%import numpy as np
%import pandas as pd
%\end{PythonBox}














\Homework
\Basic
\zhlipsum[1]
\Complex
\zhlipsum[2]




\begin{figure}
	\centering
\begin{tikzpicture}[>=stealth, thick]
	\draw (0,0)--(2,0)--(2,1)--(0,1)--(0,0);
	\fill [pattern = north east lines] (-1,-.75) rectangle (3,-.5);
	\draw(-1,-.5)--(3,-.5);
	\draw (.5,-.25) circle (.225);
	\draw (1.5,-.25) circle (.225);
	\draw (.5,-.25) circle (.1);
	\draw (1.5,-.25) circle (.1);
	\draw[->, dashed](3,0.5)--(-2,0.5);
	\draw[very thick](0,.5)--(-2,0.5);
	\foreach \i in {-.4,-.8,-1.2,-1.6,-2}
	{
		\draw (\i,0.6)--(\i,.5);
	}
	\draw (-1,1)--(-1,1.1);
	\draw (-0.6,1)--(-0.6,1.1);
	\draw[-] (-1,1)--(-0.6,1);
	\node at (-0.8,1.35){$20$N};
\end{tikzpicture}
\caption{图中的虚线表示力的作用线}
\end{figure}




\end{Project}








\Improve
\begin{QuestionItem}[2]{1}
	\item test$\Psi=\int_\Omega\left[f\left(\phi\right)+\kappa\left|\nabla\phi\right|^2\right]{\rm d}V$
	\item test
	$$\Psi=\int_\Omega\left[f\left(\phi\right)+\kappa\left|\nabla\phi\right|^2\right]{\rm d}V$$
	\item test
	\item test
	\item test
	\item test
	\item test
	\item test
\end{QuestionItem}


\Thinking
\begin{QuestionItem}{1}
	\item 测试文字
	\item 测试文字
	\item 测试文字
	\item 测试文字
\end{QuestionItem}




\chapter*{测试测试}


{
\FontSizeSet[2.0]{14pt}
测试文字

\begin{itemize}
	\item 测试文字测试文字测试文字
	\item 测试文字测试文字测试文字
	\item 测试文字测试文字测试文字
\end{itemize}

\zhlipsum
}


\chapterimage{\includegraphics[height=\paperheight]{vineyard.jpg}}
\chapter*{测试文字}
\zhlipsum[1]

\zhlipsum[2]

\zhlipsum[3]


\begin{Appendix*}
	

\chapterimage{\includegraphics[height=\paperheight]{poppies.jpg}}
\chapter{测试计数器}

\lipsum

	
\section{附录测试}



\zhlipsum[1]

\zhlipsum[2]

\section{附录测试}

\end{Appendix*}


\chapter{计数器}


\part{测试部分}


\zhlipsum
\zhlipsum


\chapterimage{\includegraphics[height=\paperheight]{garden.jpg}}
\chapter{测试计数器}
奥斯卡-克劳德·莫奈（法文：Oscar-Claude Monet，1840年11月14日——1926年12月5日），法国画家，被誉为“印象派领导者”，印象派代表人物和创始人之一。

莫奈是法国最重要的画家之一，印象派的理论和实践大部分都有他的推广。莫奈擅长光与影的实验与表现技法。他最重要的风格是改变了阴影和轮廓线的画法，在莫奈的画作中看不到非常明确的阴影，也看不到突显或平涂式的轮廓线。光和影的色彩描绘是莫奈绘画的最大特色。

从印象主义的产生、发展看，创始人并非莫奈，但真正完全实现印象主义理念和技法、并且一以贯之的当推莫奈。是他将毕生精力献给了对西方画界产生了重要影响的印象主义，是以他为首的一批艺术家的不懈努力，突破了此前学院派的保守思想，极大地冲击了19世纪后半叶占据西方画坛统治地位的官方艺术，从而为掀开西方现代绘画史新的一页，作出了重要贡献。为后人留下了宝贵的艺术财富。应该说莫奈是印象派画家中最先获得成功的人，尽管后来的野兽派、立体派、超现实主义等艺术流派，并未遵循印象派创立的一些原则，但创立这些流派的艺术家，都从印象派那里汲取过营养。\footnote{本章封面画为法国画家克劳德·莫奈的《圣阿德雷斯的阳台》。}



\section{附录测试}

\subsection{Look and Read}




\section{附录测试}

\zhlipsum
	
\cite{1992mixed}\cite{biot1941general}\cite{branquinho2004infomat}\cite{brassart2012reactive}\cite{rezaeslami2013theory}\cite{wanhassan2023physics}\cite{wanhassan2023temperature}\cite{loken2018thermal}\cite{cheng2022finite}\cite{yang2005interaction}



\begin{Appendix}
	\printbib[\subsection]{reference}*
\end{Appendix}

















